%!TEX program = xelatex
\documentclass[dvipsnames, svgnames,a4paper,11pt]{article}

\input{settings} % 导入模板的相关设置
\usepackage{lipsum}

\newcommand{\recordblank}[1]{\par\vspace*{#1}\par}

\begin{document}

\scoresTable{20分}{}{30分}{}{50分}{}{100分}{}

% \infoTable{专业}{年级}{姓名}{学号}{实验时间}{教师签名}
\infoTable{}{}{}{}{}{} % 信息留空

\begin{center}
	\LARGE 实验8 \quad 多普勒效应综合实验
\end{center}

\textbf{【实验报告注意事项】}
\begin{enumerate}
	\item 实验报告由三部分组成：
	\begin{enumerate}
		\item 预习报告：（提前一周）认真研读实验讲义，弄清实验原理；实验所需的仪器设备、用具及其使用，完成课前预习思考题；了解实验需要测量的物理量，并根据要求提前准备实验记录表格（第一循环实验已由教师提供模板，可以打印）。预习成绩低于10分（共20分）者不能做实验。
	    \item 实验记录：认真、客观记录实验条件、实验过程中的现象以及数据。实验记录请用珠笔或者钢笔书写并签名（\textcolor{red}{\textbf{用铅笔记录的被认为无效}}）。\textcolor{red}{\textbf{保持原始记录，包括写错删除部分，如因误记需要修改记录，必须按规范修改。}}（不得输入电脑打印，但可扫描手记后打印扫描件）；离开前请实验教师检查记录并签名。
	    \item 分析讨论：处理实验原始数据（学习仪器使用类型的实验除外），对数据的可靠性和合理性进行分析；按规范呈现数据和结果（图、表），包括数据、图表按顺序编号及其引用；分析物理现象（含回答实验思考题，写出问题思考过程，必要时按规范引用数据）；最后得出结论。
	\end{enumerate}
	\textbf{实验报告就是将预习报告、实验记录、和数据处理与分析合起来，加上本页封面。}
	\item 实验报告下次做实验室时交给带实验的老师（一周时间），最后一次实验结束一周后交给老师。
	\item 除实验记录外，实验报告其他部分建议双面打印。
\end{enumerate}

\textbf{【实验安全与实验室注意事项】}

\textbf{一、验证多普勒效应实验设备安装时注意事项：}
\begin{enumerate}
	\item 安装时要尽量保证红外接收器、小车上的红外发射器和超声接收器、超声发射器三者之间在同一轴线上，以保证信号传输良好；
	\item 安装时不可挤压连接电缆，以免导线折断；
	\item 安装时请确认橡胶圈是否套在主动轮上；
	\item 小车不使用时应立放，避免小车滚轮沾上污物，影响实验进行。
\end{enumerate}

\clearpage

\setcounter{section}{0}
\nsection{实验8}{多普勒效应综合实验}{预习报告}
	
\subsection{实验目的}
\begin{enumerate}
	\item 测量超声接收器运动速度与接收频率之间的关系，验证多普勒效应，并由$f$-$V$关系直线的斜率求声速。
	\item 利用多普勒效应测量物体运动过程中多个时间点的速度，查看$V$-$t$关系曲线，或调阅有关测量数据，即可得出物体在运动过程中的速度变化情况，可研究：
	\begin{enumerate}
		\item 自由落体运动，并由$V$-$t$关系直线的斜率求重力加速度。
		\item 简谐振动，可测量简谐振动的周期等参数，并与理论值比较。
		\item 匀加速直线运动，测量力、质量与加速度之间的关系，验证牛顿第二定律。
	\end{enumerate}
\end{enumerate}

\subsection{仪器用具}
\begin{table}[htbp]
	\centering
	\renewcommand\arraystretch{1.6}
	\begin{tabular}{p{0.05\textwidth}|p{0.20\textwidth}|p{0.05\textwidth}|p{0.5\textwidth}}
	\hline
	编号& 仪器用具名称 & 数量 &  主要参数（型号，测量范围，测量精度等） \\
	\hline
	1 & 多普勒效应综合实验仪 & 1 & ZKY-DPL-3 \\
	2 & 超声发射/接收器 & 1 & 谐振频率 40000-1000Hz \\
	3 & 红外发射/接收器 & 2 & \\
	4 & 导轨 & 1 & \\
	5 & 运动小车、支架、光电门、电
磁铁、弹簧、滑轮、砝码及电
机控制器等 & 1 & \\
	\hline
	\end{tabular}
\end{table}

\subsection{原理概述}

\subsubsection{一、解释超声的多普勒效应}

根据声波的多普勒效应公式，当声源与接收器之间有相对运动时，接收器接收到的频率$f$为：
\begin{equation}
	f = f_0 \frac{u - V_1\cos\alpha_1}{u - V_2\cos\alpha_2}
\end{equation}

式中$f_0$为声源发射频率，$u$为声速，$V_1$为接收器运动速率，$\alpha_1$为声源与接收器连线与接收器运动方向之间的夹角，$V_2$为声源运动速率，$\alpha_2$为声源与接收器连线与声源运动方向之间的夹角。

若声源保持不动，运动物体上的接收器沿声源与接收器连线方向以速度$V$运动，则从上式可得接收器接收到的频率应为：
\begin{equation}
	f = f_0\left(1 + \frac{V}{u}\right)
\end{equation}

当接收器向着声源运动时，$V$取正，反之取负。

若$f_0$保持不变，以光电门测量物体的运动速度，并由仪器对接收器接收到的频率自动计数，根据上式，作$f$-$V$关系图可直观验证多普勒效应，且由实验点作直线，其斜率应为$k=f_0/u$，由此可计算出声速$u=f_0/k$。

由公式(2)可解出：
\begin{equation}
	V = u\left(\frac{f}{f_0} - 1\right)
\end{equation}

若已知声速$u$及声源频率$f_0$，通过设置使仪器以某种时间间隔对接收器接收到的频率$f$采样计数，由微处理器按上式计算出接收器运动速度，由显示屏显示$V$-$t$关系图，或调阅有关测量数据，即可得出物体在运动过程中的速度变化情况，进而对物体运动状况及规律进行研究。

\subsubsection{二、简述为什么选择超声的红外调制与接收}

早期产品中，接收器接收的超声信号由导线接入实验仪进行处理。由于超声接收器安装在运动体上，导线的存在对运动状态有一定影响，导线的折断也给使用带来麻烦。

新仪器对接收到的超声信号采用了无线的红外调制-发射-接收方式。即用超声接收器信号对红外波进行调制后发射，固定在运动导轨一端的红外接收端接收红外信号后，再将超声信号解调出来。

由于红外发射/接收的过程中信号的传输是光速，远远大于声速，它引起的多普勒效应可忽略不计。采用此技术将实验中运动部分的导线去掉，使得测量更准确，操作更方便。

信号的调制-发射-接收-解调，在信号的无线传输过程中是一种常用的技术。

\subsubsection{三、简述验证多普勒效应并由测量数据计算声速的方法及步骤}

\textbf{方法：}让小车以不同速度通过光电门，仪器自动记录小车通过光电门时的平均运动速度及与之对应的平均接收频率。由仪器显示的$f$-$V$关系图可看出速度与频率的关系，若测量点成直线，符合公式(2)描述的规律，即直观验证了多普勒效应。用作图法或线性回归法计算$f$-$V$直线的斜率$k$，由$k$计算声速$u$并与声速的理论值比较，计算其百分误差。

\textbf{步骤：}
\begin{enumerate}
	\item \textbf{仪器安装：}将小车置于导轨上，使水平超声发射器、超声接收器组件（已固定在小车上）、红外接收器在同一轴线上。将组件电缆接入实验仪的对应接口上。通过连接线给小车上的传感器充电，充电时间约6～8秒。调好皮带松紧度。
	
	\item \textbf{测量准备：}实验仪开机后，输入室温$t_c$。仪器自动检测调谐频率$f_0$，将此频率记录下来。
	
	\item \textbf{测量步骤：}
	\begin{enumerate}
		\item 在液晶显示屏上，选中"多普勒效应验证实验"，并按"确认"；
		\item 修改测试总次数（选择范围5~10，一般选5次），选中"开始测试"；
		\item 用电机控制器上的"变速"按钮选定一个速度，按"确认"，再按"启动"键；
		\item 测试完成后选择"存入"或"重测"；
		\item 按电机控制器上的"变速"按钮，重新选择速度，重复上述步骤；
		\item 完成5次测量后，仪器自动存储数据，并显示$f$-$V$关系图及测量数据。
	\end{enumerate}
	
	\item \textbf{数据处理：}用作图法或线性回归法计算$f$-$V$关系直线的斜率$k$。线性回归法公式为：
	\begin{equation}
		k = \frac{\sum V_i f_i - \frac{1}{n}\sum V_i \sum f_i}{\sum V_i^2 - \frac{1}{n}(\sum V_i)^2}
	\end{equation}
	其中测量次数$n=5$。
	
	由$k$计算声速$u = f_0/k$，并与声速的理论值$u_0 = 331\sqrt{1+t_c/273}$（米/秒）比较，$t_c$表示室温（摄氏温度）。
\end{enumerate}

\subsubsection{四、简述研究自由落体运动，求自由落体加速度的方法及步骤}

\textbf{方法：}让带有超声接收器的接收组件自由下落，利用多普勒效应测量物体运动过程中多个时间点的速度，查看$V$-$t$关系曲线，并调阅有关测量数据，即可得出物体在运动过程中的速度变化情况，进而计算自由落体加速度。

\textbf{步骤：}
\begin{enumerate}
	\item \textbf{仪器安装与测量准备：}
	\begin{enumerate}
		\item 为保证超声发射器与接收器在一条垂线上，可用细绳栓住接收器组件，检查从电磁铁下垂时是否正对发射器。若对齐不好，可用底座螺钉加以调节。
		\item 充电时，让电磁阀吸住自由落体接收器组件，并让该接收器组件上充电部分和电磁阀上的九爪测试针（即充电针）接触良好。
		\item 充满电后，将接收器组件脱离充电针，下移吸附在电磁铁上。
	\end{enumerate}
	
	\item \textbf{测量步骤：}
	\begin{enumerate}
		\item 在液晶显示屏上，选中"变速运动测量实验"，并按"确认"；
		\item 修改测量点总数（选择范围8~150），选择采样步距（选择范围10~100ms），选中"开始测试"；
		\item 检查是否"失锁"，"锁定"后按"确认"按钮，电磁铁断电，接收器组件自由下落；
		\item 测量完成后，显示屏上显示$V$-$t$图，选择"数据"，阅读并记录测量结果；
		\item 可重新回到测量设置界面进行新的测量。
	\end{enumerate}
	
	\item \textbf{数据处理：}将数据记入表格中，由测量数据求得$V$-$t$直线的斜率即为重力加速度$g$。为减小偶然误差，可作多次测量，将测量的平均值作为测量值，并将测量值与理论值比较，求百分误差。考虑到断电瞬间，电磁铁可能存在剩磁，第一次采样数据的可靠性降低，故从第2个采样点开始记录数据。
\end{enumerate}

\textbf{注意事项：}
\begin{enumerate}
	\item 须将"自由落体接收器保护盒"套于发射器上，避免发射器在非正常操作时受到冲击而损坏；
	\item 安装时切不可挤压电磁阀上的电缆；
	\item 接收器组件下落时，若其运动方向不是严格的在声源与接收器的连线方向，则$\alpha$角在运动过程中增加，此时公式(2)不再严格成立，由公式(3)计算的速度误差也随之增加。故在数据处理时，可根据情况对最后2个采样点进行取舍。
\end{enumerate}

\subsubsection{五、简述研究简谐振动的方法及步骤}

\textbf{原理：}当质量为$m$的物体受到大小与位移成正比，而方向指向平衡位置的力的作用时，若以物体的运动方向为$x$轴，其运动方程为：
\begin{equation}
	m\frac{d^2x}{dt^2} = -kx
\end{equation}

由上式描述的运动称为简谐振动，当初始条件为$t = 0$时，$x = -A_0$，$V = dx/dt = 0$，则方程的解为：
\begin{equation}
	x = -A_0\cos\omega_0 t
\end{equation}

将上式对时间求导，可得速度方程：
\begin{equation}
	V = \omega_0 A_0\sin\omega_0 t
\end{equation}

由上两式可见物体作简谐振动时，位移和速度都随时间周期变化，式中$\omega_0 = \sqrt{k/m}$，为振动系统的固有角频率。

\textbf{方法：}悬挂在弹簧上的物体应作简谐振动，式中的$k$为弹簧的倔强系数。通过测量弹簧的伸长量和砝码质量可以计算$k$，通过测量速度达到最大值的时间间隔可以计算振动周期$T$和角频率$\omega$，并与理论值$\omega_0$比较。

\textbf{步骤：}
\begin{enumerate}
	\item \textbf{仪器安装与测量准备：}
	\begin{enumerate}
		\item 将弹簧悬挂于电磁铁上方的挂钩孔中，接收器组件的尾翼悬挂在弹簧上；
		\item 接收组件悬挂上弹簧之后，测量弹簧长度；
		\item 加挂质量为$m$的砝码，测量加挂砝码后弹簧的伸长量$\Delta x$，然后取下砝码；
		\item 由$m$及$\Delta x$计算$k$；
		\item 用天平称量垂直运动超声接收器组件的质量$M$，由$k$和$M$计算$\omega_0$，并与角频率的测量值$\omega$比较。
	\end{enumerate}
	
	\item \textbf{测量步骤：}
	\begin{enumerate}
		\item 在液晶显示屏上，选中"变速运动测量实验"，并按"确认"；
		\item 修改测量点总数为150（选择范围8~150），选择采样步距并修改为100（选择范围50~100ms），选中"开始测试"；
		\item 将接收器从平衡位置垂直向下拉约20cm，松手让接收器自由振荡，然后按"确认"；
		\item 接收器组件开始作简谐振动，实验仪按设置的参数自动采样；
		\item 测量完成后，显示屏上出现速度随时间变化关系的曲线。
	\end{enumerate}
	
	\item \textbf{数据处理：}查阅数据，记录第1次速度达到最大时的采样次数$N_{1\max}$和第11次速度达到最大（注：速度方向一致）时的采样次数$N_{11\max}$，就可计算实际测量的运动周期$T$及角频率$\omega$，并可计算$\omega_0$与$\omega$的百分误差。
\end{enumerate}

\subsection{实验前思考题}

\begin{question}
用多普勒效应来解释你生活中遇到的一些现象。
\end{question}
\recordblank{4cm}

\clearpage
\noindent\begin{minipage}{\textwidth}
	\renewcommand\arraystretch{1.7}
	\begin{tabularx}{\textwidth}{|X|X|X|X|}
	\hline
	专业：& \rule{2.5cm}{0.4pt} & 年级：& \rule{2.5cm}{0.4pt}\\
	\hline
	姓名：& \rule{2.5cm}{0.4pt} & 学号：& \rule{2.5cm}{0.4pt}\\
	\hline
	室温：& \rule{2.5cm}{0.4pt} & 实验地点：& \rule{2.5cm}{0.4pt}\\
	\hline
	学生签名：& \rule{2.5cm}{0.4pt} & 教师签名：& \rule{2.5cm}{0.4pt}\\
	\hline
	实验时间：& \rule{2.5cm}{0.4pt} & 教师签名：& \rule{2.5cm}{0.4pt}\\
	\hline
	\end{tabularx}
\vspace{0.8em}
\refstepcounter{section}
\phantomsection
\addcontentsline{toc}{section}{实验8 多普勒效应综合实验 实验记录}
\begin{center}
	\LARGE 实验8 多普勒效应综合实验 \hspace{11pt} \textbf{实验记录}
\end{center}
\subsection{一、验证多普勒效应并由测量数据计算声速}

\subsubsection{1. 实验数据记录}
\recordblank{5.0cm}
\end{minipage}

\subsubsection{2. 根据原始测量数据，做出$f$-$V$曲线。计算直线斜率$k$ (1/m)和声速测量值$u=f_0/k$ (m/s)。将计算值与声速理论值进行比较，分析误差。}
\recordblank{5.0cm}
\subsection{二、研究自由落体运动，求自由落体加速度}

\subsubsection{1. 实验数据记录}
\recordblank{9.8cm}
\subsubsection{2. 根据上表原始测量数据，计算重力加速度的值。将实验计算值与当地的重力加速度理论值进行比较，并分析误差。}
\recordblank{5.0cm}
\subsection{三、研究简谐振动}

\subsubsection{1. 实验数据记录}
\recordblank{5.0cm}
\subsubsection{2. 根据原始测量数据，计算弹簧的振动角频率$\omega=2\pi/T$；根据弹簧的参数，计算固有角频率；将实验测得的角频率数据与理论值比较，分析误差。}
\recordblank{8.4cm}
\subsection{四、研究匀变速直线运动，验证牛顿第二运动定律（选做题）}

\subsubsection{1. 实验数据记录}
\recordblank{8.4cm}
\subsubsection{2. 以实验得出的加速度$a$作纵轴，$[(1-C)M-m]/[(1-C)M+m+J/R^2]$作横轴作图。分析上图中的关系是否为线性关系？若是线性关系，求出直线的斜率即重力加速度。说明符合实验讲义中(9)式描述的规律，验证了牛顿第二定律。}
\recordblank{6.2cm}
\section{附录：实验原始记录}


\recordblank{7cm}


\end{document}
